\section{Degree sums, common trimming, and normalization}
\label{l:foundations}

We first supply the three ingredients needed to complete a transfer: a uniform count of small representation degrees, one measure satisfying all coset bounds at once, and the precise Fourier conversion to total variation. For a finite set, $U$ denotes uniform probability. For a subgroup $H\le G$, the pushforward of $\mu$ to $G/H$ is denoted by $\mu_H$, so $\mu_H(gH)=\mu(gH)$. Our convention is $\|\mu-\nu\|_{\TV}=\frac12\sum_x|\mu(x)-\nu(x)|$ for probabilities.

\begin{lemma}[A short reciprocal-degree proof]\label{l:degree-basic}
For every $u>0$, $C_u<\infty$ and
\[
 Z_u(n):=\sum_{\lambda\vdash n,\ \lambda\notin\{(n),(1^n)\}}D_\lambda^{-u}\longrightarrow0.
\]
The excluded partitions form a set, so the unique partition of one is removed only once.
\end{lemma}
\begin{proof}
We use the classical tableau dimension and hook formulas: $D_\lambda$ is the number of fillings of the diagram of $\lambda$ by $1,\ldots,n$ increasing along each row and column, and
\[D_\lambda=\frac{n!}{\prod_{x\in\lambda}h(x)},\]
where $h(x)$ counts the box $x$ and the boxes strictly right of it or strictly below it \cite{frame-robinson-thrall,James1978}. Write $p(k)$ for the number of partitions of $k$, with $p(0)=1$. Euler's product at $x=e^{-1/\sqrt{k}}$ gives, for $k\ge1$,
\[
 \log p(k)\le\sqrt{k}+\sum_{j\ge1}\frac1{j(e^{j/\sqrt{k}}-1)}
 \le\sqrt{k}\left(1+\sum_{j\ge1}j^{-2}\right)\le3\sqrt{k}.
\]
If a Young diagram has a top row of length $a$ and $j\le a$ lower boxes, its number of standard tableaux is at least $\binom aj$. Put the first $j$ labels in the first $j$ top-row positions, choose the $j$ lower labels from the remaining $a$ labels, and fill the lower diagram in one fixed standard relative order. Its width is at most $j$, so every top-row predecessor is smaller. Fill the rest of the row increasingly. A tableau of any Young subdiagram extends to the full diagram by adding larger labels in an order respecting row and column predecessors.

Orient a longest row or column as the first row and let its length be $n-j$. When $1\le j\le n/3$, at most $2p(j)$ diagrams have this value and
\[
 D_\lambda\ge\binom{n-j}{j}\ge2^j.
\]
For each fixed $j>0$ the binomial coefficient diverges, and $2p(j)2^{-uj}$ is summable. Dominated convergence handles this range.

In the remaining range put $a=\max(\lambda_1,\lambda'_1)<2n/3$. If $a\ge n/10$, retain the first row and a lower Young order ideal of $\lfloor a/2\rfloor$ boxes; there are enough lower boxes. Extension and the preceding construction give $D_\lambda\ge2^a/(a+1)$, exponential in $n$. If $a<n/10$, every hook has length at most $2a$, and the hook formula gives $D_\lambda\ge n!/(2a)^n\ge(5/e)^n$. Both bounds beat $p(n)\le e^{3\sqrt n}$. Thus $Z_u(n)\to0$. Adding the row and column shows that the full sum tends to two; finitely many smaller values are finite.
\end{proof}

\begin{lemma}[Selecting one partition]\label{l:partition-selection}
If $n$ is divisible by $b$ and a law $\mu$ on $S_n$ satisfies
\[
 \E_{|M|=n/b}\|\mu_{H_M}-U_{G/H_M}\|_{\TV}\le\delta,
\]
there is a partition into $b$ equal blocks for which the sum of the complementary coset errors is at most $b\delta$.
\end{lemma}
\begin{proof}
Every block of a uniformly chosen ordered equal partition is itself a uniform $n/b$-subset. The expected sum of the errors is at most $b\delta$, so one partition has sum no larger than that expectation.
\end{proof}

\begin{lemma}[Two common trimming operations]\label{l:trim}
Let $H_1,\ldots,H_b$ be subgroups of a finite group $G$, let $\mu$ be a probability, and put $\delta_i=\|\mu_{H_i}-U_{G/H_i}\|_{\TV}$ and $\delta=\sum_i\delta_i$.
There is a subprobability $f\le\mu$ with
\[
 f(G)\ge1-\delta,\qquad f(gH_i)\le[G:H_i]^{-1}\quad(g\in G,i\le b).
\]
For any $c>1$, deleting the union of all cosets with original density greater than $c$ loses mass at most $c\delta/(c-1)$. If that loss is $\eta<1$, the conditional law $\nu$ has
\[
 \|\nu-\mu\|_{\TV}=\eta,\qquad \nu(gH_i)\le\frac{c}{1-\eta}[G:H_i]^{-1}.
\]
In particular $c=2$ loses at most $2\delta$.
\end{lemma}
\begin{proof}
Process the subgroup coset systems successively. Within each coset whose current mass exceeds its uniform mass, multiply all current point masses by the ratio of the cap to the current mass. The loss at stage $i$ is at most the original positive excess on the $i$th quotient, namely $\delta_i$, because all current masses are below $\mu$. Later decreases preserve earlier caps.

There is also a simultaneous construction with the same guarantee. For each $i$, trim $\mu$ separately to a measure $f_i$, and set $f(g)=\min_i f_i(g)$. Then $f\le f_i$ ensures every cap, while $\mu-f\le\sum_i(\mu-f_i)$ bounds the loss by $\delta$. Thus these two constructions need not produce the same measure, but prove exactly the same assertion.

If a coset $Q$ has $\mu(Q)>cU(Q)$, then $\mu(Q)\le c(\mu(Q)-U(Q))/(c-1)$. Summing over its heavy cosets and then over subgroup systems bounds the removed mass. A surviving coset has original mass at most $cU(Q)$; division by $1-\eta$ gives the displayed cap. Conditioning deletes mass $\eta$ and adds that same mass on the retained set, proving the variation equality.
\end{proof}

\begin{lemma}[Plancherel with a mass deficit]\label{l:plancherel}
For every complex mass function $a$ on a finite group $G$,
\[
 |G|\sum_g|a(g)|^2=\sum_{\rho\in\Irr(G)}D_\rho\|\widehat a(\rho)\|_{\HS}^2.
\]
If $f\ge0$ has mass $z\le1$, then
\begin{align}
 |G|\sum_g|f(g)-zU(g)|^2&=\sum_{\rho\ne\one}D_\rho\|\widehat f(\rho)\|_{\HS}^2,\label{l:centered-plancherel}\\
 \left(\sum_g|f(g)-U(g)|\right)^2&\le(1-z)^2+\sum_{\rho\ne\one}D_\rho\|\widehat f(\rho)\|_{\HS}^2.\label{l:deficit-plancherel}
\end{align}
If $f\le\mu$ for a probability $\mu$, its normalization, when $z>0$, is at variation distance at most $1-z$ from $\mu$.
\end{lemma}
\begin{proof}
Expand the squared Hilbert--Schmidt norms. The regular-character identity
\[\sum_\rho D_\rho\chi_\rho(h^{-1}g)=|G|\one_{g=h}\]
proves Plancherel. Apply it to $f-zU$, whose trivial coefficient is zero, and to $f-U$, whose trivial coefficient is $z-1$. Cauchy--Schwarz gives \eqref{l:deficit-plancherel}. Finally, $\sum|f/z-f|=1-z=\sum|\mu-f|$, so the triangle inequality proves the normalization claim.
\end{proof}

With convolution $(f*k)(g)=\sum_xf(x)k(x^{-1}g)$, transforms multiply in the displayed order. If $f\le\mu$, positivity implies $f^{*r}\le\mu^{*r}$ and the missing mass is $1-z^r\le r(1-z)$. If $\mu$ has sign mean $s$, then $|\widehat f(\sgn)|\le |s|+1-z$. For nearby probabilities, the corresponding sign perturbation is at most twice their total-variation distance.

\begin{proposition}[Fourier amplification]\label{l:amplification}
Let $\mu_n$ be probabilities on $S_n$ with $\widehat\mu_n(\sgn)=o(1)$. Suppose either $f_n$ is a nonnegative subprobability with $f_n\le\mu_n$ and mass tending to one, or $f_n$ is a probability with $\|f_n-\mu_n\|_{\TV}=o(1)$. If, uniformly for every irreducible of degree $D>1$,
\[
 \|\widehat f_n\|_{\op}\le A D^{-\beta},
\]
where $A,\beta>0$ are fixed, then $\|\mu_n^{*r}-U\|_{\TV}\to0$ for every fixed integer $r$ with $2-2r\beta<0$.
If instead $\|\widehat f_n\|_{\HS}^2\le A D^{-\gamma}$, the sufficient condition is $1-r\gamma<0$.
\end{proposition}
\begin{proof}
The sign term tends to zero by the preceding mass bounds. In the operator case,
\[
 D\|\widehat f_n^{\,r}\|_{\HS}^2\le D^2\|\widehat f_n\|_{\op}^{2r}\le A^{2r}D^{2-2r\beta}.
\]
In the Hilbert--Schmidt case submultiplicativity gives $D\|\widehat f_n^{\,r}\|_{\HS}^2\le A^rD^{1-r\gamma}$. The sums tend to zero by Lemma~\ref{l:degree-basic}. The trivial coefficient tends to one. Lemma~\ref{l:plancherel} gives convergence for $f_n^{*r}$; the missing mass or a telescoping replacement of the $r$ probability factors gives convergence for $\mu_n^{*r}$. No normality or diagonalization is used.
\end{proof}
